%! Characteristics of Radical Functions — Unit 6, Lesson 6.0 — Warm-Up (Answer Key)
\documentclass[10pt]{article}
\usepackage{algebra2-article}
\usepackage{algebra2-key}   % pulls in -boxes; do NOT also load -boxes
\newcommand{\TallMath}[1]{$\displaystyle #1\rule[-1.4em]{0pt}{3.2em}$}

\begin{document}

\pageheader{Unit 6, Lesson 6.0}{Warm-Up --- Answer Key}
\namedateperiod

\vspace{0.10in}

\begin{notesbox}{Getting Ready: Skills We'll Use Today}
\small These three skills power today's lesson on reading radical graphs. Work quickly.

\vspace{4pt}
\textbf{1. Which roots actually exist?} Evaluate each one. If there is no real answer, write
``not a real number.''
\begin{center}
\renewcommand{\arraystretch}{1.5}
\begin{tabular}{@{}l l@{}}
(a) \; $\sqrt{36}=$ \ans{6} & (b) \; $\sqrt{-36}=$ \ans{not a real number} \\
(c) \; $\sqrt[3]{8}=$ \ans{2} & (d) \; $\sqrt[3]{-8}=$ \ans{$-2$} \\
\end{tabular}
\end{center}
Exactly one of the four has no real answer. Which one --- and what is it about \emph{squaring} that
makes it impossible? \ansline{(b). No real number squares to a negative --- a positive squared is
positive and a negative squared is also positive --- so nothing squares to $-36$. Cubing keeps the
sign, so $\sqrt[3]{-8}$ is fine.}

\vspace{4pt}
\textbf{2. Solve each inequality.} (Watch the last one: dividing or multiplying by a negative
\emph{flips} the sign.)
\begin{enumerate}[label=(\alph*), leftmargin=2.1em, itemsep=3pt]
  \item $x-3\ge 0$ \quad$\Rightarrow$\quad \ans{$x\ge 3$}
  \item $2x+8\ge 0$ \quad$\Rightarrow$\quad \ans{$x\ge -4$}
  \item $5-x\ge 0$ \quad$\Rightarrow$\quad \ans{$x\le 5$}
\end{enumerate}

\vspace{4pt}
\textbf{3. Two tables that are hiding something.} Complete both.
\begin{center}
\renewcommand{\arraystretch}{1.3}
\begin{tabular}{|c|c|c|c|c|c|}
\hline
\multicolumn{6}{|c|}{\small $y=x^2$} \\ \hline
$x$ & $0$ & $1$ & $2$ & $3$ & $4$ \\ \hline
$y$ & \ans{0} & \ans{1} & \ans{4} & \ans{9} & \ans{16} \\ \hline
\end{tabular}
\hspace{0.2in}
\begin{tabular}{|c|c|c|c|c|c|}
\hline
\multicolumn{6}{|c|}{\small $y=\sqrt{x}$} \\ \hline
$x$ & $0$ & $1$ & $4$ & $9$ & $16$ \\ \hline
$y$ & \ans{0} & \ans{1} & \ans{2} & \ans{3} & \ans{4} \\ \hline
\end{tabular}
\end{center}
Compare the two tables. What happened to the $x$-row and the $y$-row? \ansline{They traded places ---
the inputs of one table are the outputs of the other. Squaring and taking a square root undo each
other.}

\end{notesbox}

\begin{teachernote}
Each item seeds one of today's new ideas, so debrief all three aloud.
\textbf{Item 1} is the whole domain story in miniature: an \emph{even} index refuses a negative
radicand, an \emph{odd} index accepts one. That single fact is why $y=\sqrt{x}$ is only half a graph
while $y=\sqrt[3]{x}$ runs across the entire plane --- do not give away that consequence yet.
\textbf{Item 2} is the domain-restriction engine, reused all unit: to find a square root function's
domain, students will set the \emph{radicand} $\ge 0$ and solve, and part~(c) plants the sign-flip
that produces a graph opening to the \emph{left}. \textbf{Item 3} plants the inverse relationship
without naming it; the swapped rows are exactly the reflection over $y=x$ that Section~4 of the notes
makes visible, and the fact that the $y=x^2$ table was started at $x=0$ (not $x=-4$) is the seed of
the restricted domain that Lesson~6.7 will need. Do \emph{not} say ``index,'' ``radicand,'' or
``inverse'' yet --- let the Hook graphs earn those words.
\end{teachernote}

\end{document}
